\documentclass[10pt]{beamer}
\usetheme{metropolis}
\usepackage{graphicx}
\usepackage{listings}
\usepackage{booktabs}
\usepackage{xcolor}
\definecolor{codebg}{HTML}{F5F5F5}
\definecolor{codegreen}{HTML}{2E7D32}
\definecolor{codegray}{HTML}{757575}
\definecolor{codeblue}{HTML}{1565C0}
\lstdefinestyle{rstyle}{language=R,backgroundcolor=\color{codebg},basicstyle=\ttfamily\scriptsize,
  keywordstyle=\color{codeblue}\bfseries,stringstyle=\color{codegreen},
  commentstyle=\color{codegray}\itshape,breaklines=true,frame=single,
  rulecolor=\color{codegray!30},showstringspaces=false,tabsize=2,
  morekeywords={library,pheatmap,ComplexHeatmap,Heatmap,HeatmapAnnotation,
    hclust,dist,scale,cor}}
\lstdefinestyle{pythonstyle}{language=Python,backgroundcolor=\color{codebg},basicstyle=\ttfamily\scriptsize,
  keywordstyle=\color{codeblue}\bfseries,stringstyle=\color{codegreen},
  commentstyle=\color{codegray}\itshape,breaklines=true,frame=single,
  rulecolor=\color{codegray!30},showstringspaces=false,tabsize=4,
  morekeywords={import,from,as,pd,np,plt,sns,clustermap,heatmap,
    linkage,dendrogram,zscore,StandardScaler}}
\title{Heatmaps \& Clustered Displays}
\subtitle{Workshop 5 --- Module 4}
\author{Prof.Asoc.\ Endri Raco}
\institute{Data Visualization for Data Scientists \\ UNYT Tirana}
\date{2025/26}
\begin{document}

\begin{frame}
  \titlepage
\end{frame}

\begin{frame}{Module Roadmap}
  \begin{enumerate}
    \item The heatmap as a multivariate display
    \item Clustered heatmaps: dendrograms reorder rows and columns
    \item Z-score normalisation: revealing within-row patterns
    \item Five design principles for publication-quality heatmaps
    \item Packages: pheatmap, ComplexHeatmap, sns.clustermap
  \end{enumerate}
  \begin{exampleblock}{Learning Outcome}
    You will produce clustered heatmaps with dendrograms and annotation
    tracks, apply z-score normalisation to reveal feature patterns, and
    choose between sequential and diverging palettes based on the data
    semantics.
  \end{exampleblock}
\end{frame}

\section{The Heatmap}

\begin{frame}{Rows $\times$ Columns, Colour = Value}
  \begin{center}
    \includegraphics[width=0.72\textwidth]{figures_w05m04/basic_heatmap}
  \end{center}
  \begin{exampleblock}{Data Insight}
    A heatmap encodes a matrix of values as coloured cells. Rows
    = observations (samples, genes, customers). Columns = variables
    (features, time points, metrics). Colour intensity = the value.
    Without clustering, patterns are hidden in the row/column order.
  \end{exampleblock}
\end{frame}

\section{Clustered Heatmaps}

\begin{frame}{Dendrograms Reveal Block Structure}
  \begin{center}
    \includegraphics[width=0.72\textwidth]{figures_w05m04/clustered_heatmap}
  \end{center}
  \begin{alertblock}{Pro-Tip}
    Hierarchical clustering reorders both rows and columns so that
    similar profiles are adjacent. The resulting \textbf{block patterns}
    reveal which features co-vary and which samples share profiles.
    Without clustering, the same data looks like noise.
  \end{alertblock}
\end{frame}

\begin{frame}[fragile]{Clustered Heatmap in Code}
  \begin{columns}[T]
    \begin{column}{0.48\textwidth}
      \textbf{R (pheatmap)}
\begin{lstlisting}[style=rstyle]
library(pheatmap)

pheatmap(mat,
  scale = "row",          # z-score
  clustering_method = "ward.D2",
  clustering_distance_rows =
    "euclidean",
  color = colorRampPalette(
    c("#1565C0","white","#E53935"))(100),
  border_color = NA,
  fontsize_row = 6,
  fontsize_col = 8,
  annotation_row = ann_df, # metadata
  main = "Clustered Heatmap")

# ComplexHeatmap (Bioconductor)
library(ComplexHeatmap)
Heatmap(mat,
  name = "z-score",
  row_split = 3,      # k-means split
  column_split = 2,
  top_annotation =
    HeatmapAnnotation(type = ann))
\end{lstlisting}
    \end{column}
    \begin{column}{0.48\textwidth}
      \textbf{Python (seaborn)}
\begin{lstlisting}[style=pythonstyle]
import seaborn as sns

g = sns.clustermap(
    df,
    method="ward",
    metric="euclidean",
    standard_scale=0,    # z-score rows
    cmap="RdBu_r",
    figsize=(8, 10),
    linewidths=0.5,
    col_cluster=True,
    row_cluster=True,
    cbar_kws={"label": "z-score"},
    # Annotation sidebar
    row_colors=cluster_colors,
    col_colors=feature_colors)

g.ax_heatmap.set_xticklabels(
    g.ax_heatmap.get_xticklabels(),
    rotation=45, fontsize=7)
\end{lstlisting}
    \end{column}
  \end{columns}
\end{frame}

\section{Z-Score Normalisation}

\begin{frame}{Raw vs Z-Scored: Patterns Emerge}
  \begin{center}
    \includegraphics[width=0.88\textwidth]{figures_w05m04/zscore_effect}
  \end{center}
\end{frame}

\begin{frame}{When to Z-Score}
  \centering
  \begin{tabular}{lll}
    \toprule
    \textbf{Situation} & \textbf{Normalise?} & \textbf{Palette} \\
    \midrule
    Variables on different scales & Yes (z-score rows) & Diverging (RdBu) \\
    Absolute values matter & No & Sequential (YlOrRd) \\
    Correlation matrix & Already $[-1,1]$ & Diverging (RdBu) \\
    Count matrix (e.g., word freq) & Log-transform first & Sequential (viridis) \\
    Comparing relative profiles & Yes (z-score rows) & Diverging \\
    \bottomrule
  \end{tabular}
  \vspace{6pt}
  {\small Z-scoring per row: $z_{ij} = (x_{ij} - \bar{x}_i) / s_i$.
  Centres each row at 0 with unit variance, so colour encodes
  ``how many SDs above/below this row's mean'' rather than
  the raw value.}
\end{frame}

\section{Design Principles}

\begin{frame}{Five Heatmap Design Principles}
  \begin{center}
    \includegraphics[width=0.92\textwidth]{figures_w05m04/design}
  \end{center}
\end{frame}

\section{Packages}

\begin{frame}{Heatmap Packages}
  \begin{center}
    \includegraphics[width=0.88\textwidth]{figures_w05m04/packages}
  \end{center}
\end{frame}

\begin{frame}{When to Use Heatmaps vs Other Charts}
  \centering
  \begin{tabular}{lll}
    \toprule
    \textbf{Data} & \textbf{Best Chart} & \textbf{Why Not Heatmap} \\
    \midrule
    5--50 observations $\times$ 5--50 features & Heatmap & \emph{This is the sweet spot} \\
    2 variables & Scatter & Heatmap wastes 1D \\
    1 variable over time & Line chart & Temporal axis natural \\
    $>$500 rows & Heatmap on summary stats & Too many rows unreadable \\
    Correlation matrix & Heatmap (triangle) & \emph{Ideal use case} \\
    Geographic data & Map & Spatial layout essential \\
    \bottomrule
  \end{tabular}
\end{frame}

\section{Synthesis}

\begin{frame}{Key Takeaways}
  \begin{enumerate}
    \item Heatmaps encode a \textbf{matrix} as coloured cells.
          Their power comes from \textbf{clustering}: reorder rows
          and columns to reveal block patterns.
    \item \textbf{Z-score rows} when variables have different scales;
          use \textbf{diverging palette} centred at 0.
    \item \textbf{Sequential palette} (YlOrRd, viridis) for magnitude;
          \textbf{diverging} (RdBu) for deviation from a reference.
    \item In R: \texttt{pheatmap} (quick), \texttt{ComplexHeatmap}
          (publication-quality, annotation tracks).
          In Python: \texttt{sns.clustermap} (auto dendrograms).
    \item Sweet spot: 5--50 rows $\times$ 5--50 columns. Beyond that,
          aggregate or subset first.
  \end{enumerate}
\end{frame}

\begin{frame}{Essential Readings}
  \begin{itemize}
    \item Wilkinson, L. \& Friendly, M. (2009). ``The History of the
          Cluster Heat Map.'' \emph{The American Statistician}, 63(2).
    \item Gu, Z. et al. (2016). ``ComplexHeatmap.''
          \emph{Bioinformatics}, 32(18).
    \item Wilke, C.~O. (2019). \emph{Fundamentals of Data Visualization},
          Chapter 6.
  \end{itemize}
  \vspace{6pt}
  \textbf{Next Module}: Introduction to Spatial Data (W05-M05)
\end{frame}

\end{document}
